\chapter{Introduction}
\label{ch:introduction}

\section{Why Formal Methods Stay Niche}

Formal methods --- the use of precise, machine-checkable specifications to prove
that software behaves as intended --- have existed for half a century, rest on
solid theory, and are demonstrably effective where they are applied. And yet, by
any honest measure, they remain confined to a few niche corners of software
engineering: safety-critical avionics and rail, security-critical kernels and
cryptographic code, the occasional high-assurance compiler. The overwhelming
majority of software --- the libraries, the services, the applications that make
up everyday practice --- is built and shipped without a formal specification in
sight. This is the puzzle the thesis begins from. The technology to verify
software against a specification is mature and freely available; the benefits of
doing so are not in dispute; and still the practice does not spread. Why?

The thesis's answer, which it then sets out to act on, is that the binding
constraint is not the verification technology but the \emph{specifications it
requires as input}. A verifier checks code against a specification, and someone
must write that specification. Writing it by hand is expensive, the expense
scales with the size of the codebase, and it is paid up front before any benefit
is realised --- a combination that makes formal verification economically
rational only where the cost of failure is catastrophic, which is precisely the
niche where it is used. If the specifications could be produced automatically,
the economics would change, and the question of this thesis is whether automatic
inference --- and the capabilities inferred specifications enable --- can move
formal methods out of the niche. The remainder of this section develops why the
specification is the bottleneck; the sections that follow set out the two ways
the thesis argues automatic inference attacks it.

\subsection{What Code Does Versus What It Should Do}

Software engineering has, since its inception, been organised around a
distinction that is easy to state and hard to operationalise: the difference
between what a program \emph{does} and what it is \emph{intended to do}. The
gap between the two is where bugs live. A test that checks only what the code
does will pass on buggy code; documentation that describes only what the code
does drifts the moment the code changes; a developer reading an unfamiliar
library learns what its methods do by reading their bodies, but learns what
they are \emph{for} only by inference, guesswork, or asking someone else. The
intended behaviour of a method is a real and consequential property, and yet
it is rarely written down in a form a machine can check.

The classical answer to this problem is the \emph{formal
specification}~\cite{meyer1992design, jml, leavens2006design}: a precise,
machine-checkable statement of a program's intended behaviour, expressed as
preconditions (what must hold when a method is called), postconditions (what
the method guarantees on return), and invariants (what remains true
throughout). A specification is the externalisation of intent. Where a test
encodes intent implicitly and by example, a specification encodes it
explicitly and universally; where documentation states intent in prose that
no tool can verify, a specification states it in logic that a verifier can
discharge. The benefits of having specifications are not in dispute. They
solve the test-oracle problem by supplying an independent statement of
correct behaviour against which generated tests can assert. They solve the
documentation-drift problem by being checkable, and therefore falsifiable,
against the code they describe. They make it possible to reason modularly
about a system, replacing a called method's body with its contract. They are,
in short, the substrate on which formal verification, contract-based design,
specification-based testing, and a great deal of static analysis are built.

The benefits are not in dispute; the cost is. Writing specifications by hand
is expensive, and the expense is not incidental but structural. In
large-scale formal-verification efforts the specification burden has been
measured at roughly half of total project
effort~\cite{andronick_large-scale_2012, klein2014sel4, matichuk_cost_2015}.
The cost is paid up front, before any verification benefit is realised, and
it scales with the size of the codebase rather than with the difficulty of
the bugs one hopes to catch. Surveys of why formal methods have not achieved
broad industrial adoption return repeatedly to this
point~\cite{hall1990seven, zimmermann2019lightweight}: the technology to
\emph{check} specifications has matured to the point of practicality, but the
labour to \emph{write} them has not been removed, and so the technology sits
unused. The barrier is not the verifier. The barrier is the blank page.

This thesis takes as its premise that the blank page is the problem worth
attacking, and that the way to attack it is to infer specifications
automatically from the code that already exists. If a static analyser can
read a method body and emit a plausible, formally checkable contract --- a
\texttt{requires} clause for each guard, an \texttt{ensures} clause for each
return path, a loop invariant for each loop --- then the up-front labour cost
collapses, and the question changes from ``can we afford to write
specifications?'' to ``what becomes possible once we have them for free?''
That second question is the one this thesis pursues.

\section{The Specification Gap in Practice}

The abstract argument above has concrete manifestations in everyday software
practice, and three of them recur throughout this thesis as motivating
scenarios.

The first is the unit-testing scenario. A developer tasked with testing a method
faces the oracle problem in its purest form: the test framework will run any
assertion the developer writes, but it cannot tell the developer what the
assertion should be. In the absence of a statement of intent, the path of least
resistance is to assert what the code currently produces --- to run the method,
observe the output, and bake that output into the test as the expected value.
The resulting test is a regression test, useful for detecting future change but
useless for detecting present error, because it was written by consulting the
very implementation it purports to check. When the testing is delegated to a
language model, the same dynamic plays out automatically and at scale: the model,
given the body, asserts what the body does. A specification breaks this loop by
supplying an independent statement of intent, and the first study of this thesis
measures exactly how much that matters.

The second is the integration scenario. A developer adopting a third-party
library receives it as a compiled artefact --- a JAR of bytecode, a remote
service behind an HTTP interface --- and must learn how to call it correctly.
The compiled artefact carries the method signatures, which give the types, and
perhaps some documentation, which gives prose; neither gives a checkable
statement of the preconditions the library's methods impose or the postconditions
they guarantee. The developer discovers the contract by trial, by reading
whatever source is available, or by encountering the exceptions that a violated
precondition throws at runtime. A specification carried inside the artefact would
let tooling --- an IDE, a verifier, a contract-checking test generator --- enforce
the contract at the call site, but only if the specification can travel with the
compiled form, which is the problem the second study addresses.

The third is the reasoning scenario, increasingly mediated by language models. A
model asked to use, extend, or generate code against a library reasons about that
library from whatever it can see: the signatures, the documentation, and whatever
it absorbed of the library during training. For a library that postdates the
model's training cut-off, or a private library it never saw, the model has only
the signatures and the documentation, and it reasons about the library's
behaviour by guesswork. A machine-readable specification, supplied as context,
replaces the guesswork with a checkable contract --- and a model is, as the first
study shows, an able consumer of such contracts. The same machinery that helps a
model write better tests can in principle help it write better client code,
which is one of the broader motivations for making inferred specifications both
abundant and distributable.

The fourth, and for this thesis the most consequential, is the
\emph{library-evolution} scenario. A library is a moving target: it is released,
revised, and re-released, and each new version may change the behaviour its
methods promise --- tightening a precondition, weakening a guarantee, altering a
corner case. A consumer who upgrades a dependency wants to know whether the new
version is compatible with their code, and today this question is answered
crudely. Semantic versioning conventions signal intent through a version number;
signature compatibility is checked by the compiler; but \emph{behavioural}
compatibility --- whether the new version still honours the contracts the old one
did, and whether the consumer's code still satisfies what the new version
requires --- is left to tests, to release notes, and to discovery in production.
A specification at the library boundary changes this. Because a verifier checks a
caller against its callees' contracts, and because a caller's own inferred
specification is \emph{derived from} those callee contracts, a change in a
library's specifications propagates into the specifications of the code that uses
it. When the library's inferred contracts change between versions, the impact on
a consumer's contracts becomes computable, and behavioural compatibility becomes a
question answerable at the specification level rather than guessed. This scenario
is the one the thesis's second and third studies build toward, and it is the
clearest instance of inferred specifications enabling something that manual
specification, confined to the systems that can afford it, never made routine.

These scenarios share a structure. In each, a consumer --- a developer, a test
generator, a model, a team upgrading a dependency --- needs to know what a piece
of code is intended to do, and in each the intent is present in the code's design
but absent from the form in which the consumer encounters it.

\section{Two Routes Out of the Niche}

If the manual cost of specifications is what confines formal methods to the
niche, then automatic inference attacks the confinement along two routes, and
the body of this thesis is organised around them.

The first route is the obvious one: \emph{remove the cost}. If a static analyser
can read a method and emit a formally checkable contract, the up-front labour
that made verification economically rational only for catastrophe-cost systems
collapses, and the question changes from ``can we afford to write
specifications?'' to ``what becomes possible once we have them cheaply?''. This
route raises two questions the thesis must answer before the second route is even
worth pursuing. The first is whether inference is \emph{accurate} enough: a
specification inferred by heuristic pattern matching is not guaranteed sound, so
how well do the inferred contracts match both documented developer intent and the
verdict of a formal verifier? The second is whether inferred specifications are
\emph{worth having}: granting they can be produced cheaply, do they deliver
measurable value to a consumer, or is a cheap specification a worthless one? The
first study (Chapters~\ref{ch:inference} and~\ref{ch:testgen}) answers both,
using a large-language-model test generator --- the most demanding contemporary
consumer of code-level intent, and one that suffers acutely from the oracle
problem~\cite{konstantinou2024oracles, schafer2023adaptive} --- as the instrument
that measures the value.

The second route is the one that makes inferred specifications more than a
cheaper version of what experts already write, and it is where the thesis's
distinctive contribution lies: \emph{make specifications a property of the
library boundary, so that verification and compatibility analysis become routine
across the dependencies that make up real software}. This route depends on two
facts about how verification and inference compose. First, a verifier checks a
caller against its callees' \emph{contracts}, not their bodies, so if a
dependency's contract is available a consumer can verify client code against it
without the dependency's source --- which requires that the contract travel with
the compiled library. Second, and more consequentially, a caller's own inferred
specification is \emph{derived from} its callees' specifications, so a change to a
library propagates into the contracts of the code that uses it: when a library's
inferred specifications change between versions, the effect on a consumer's
contracts --- and hence the behavioural compatibility of the upgrade --- becomes
computable. The second and third studies (Chapters~\ref{ch:distribution}
and~\ref{ch:compositional}) establish the two facts this route rests on: that
inferred specifications can be carried losslessly with compiled libraries and
service interfaces, and that caller specifications depend substantially and
measurably on callee specifications, through exactly the cross-method obligations
that a single-pass analysis cannot see and a compositional one recovers.

The two routes are not independent. The cost route produces the specifications;
the boundary route makes them useful across the dependency graph that real
software is built from. And the boundary route is what answers the adoption
puzzle most directly: a capability that did not exist before --- behavioural
compatibility analysis at the library level, available without a dependency's
source --- is the kind of new value that moves a technology out of its niche,
because it offers practitioners something they want and could not previously
have, rather than a cheaper version of something most of them never used.

\section{Thesis Statement}

The central claim of this thesis can be stated in one sentence:

\begin{quote}
\emph{The adoption of formal methods is constrained chiefly by the manual cost of
the specifications they require, and that constraint is substantially removable:
heuristic static inference, validated against a formal oracle, can produce Java
Modeling Language specifications accurately and cheaply; those specifications can
be carried with the compiled libraries and services that consume them; and
because a caller's specification is derived from its callees', specifications at
the library boundary make verification and behavioural-compatibility analysis ---
capabilities formerly confined to systems that could afford hand-written
specifications --- available across ordinary software.}
\end{quote}

The claim has two halves, corresponding to the two routes of the previous section.
The first half --- that inference removes the cost barrier --- is the necessary
groundwork: it must be true, and demonstrably so, before the second half is worth
arguing. The second half --- that library-boundary specifications enable
verification across dependencies and, through the dependence of caller
specifications on callee specifications, behavioural-compatibility analysis --- is
the half that bears on adoption most directly, because it offers a capability
manual specification never made routine rather than a cheaper version of one it
did. The thesis defends the first half with a fidelity measurement and a
controlled value experiment, and it establishes the mechanisms the second half
rests on --- lossless distribution, and the substantial dependence of caller
specifications on callee specifications --- by measurement on real libraries,
positioning the full compatibility application as the capability these mechanisms
enable.

The claim is deliberately empirical. It does not assert that heuristic
inference is sound in the sense a logician would demand; it asserts that
inference plus \emph{post-hoc} verification yields specifications whose accuracy
can be measured, whose downstream value can be demonstrated, and whose
cross-library propagation can be quantified. Each part is defended with
quantitative evidence on real-world code rather than with appeal to construction.

\section{Research Questions}

The thesis is structured around four research questions. The first two ask
whether inference removes the cost barrier; the second two establish the
mechanisms by which library-boundary specifications enable verification and
compatibility analysis.

\begin{description}
  \item[RQ1 (Inference accuracy).] How accurately can JML preconditions,
        postconditions, and loop invariants be inferred from Java source by
        weakest-precondition/strongest-postcondition-organised pattern matching
        --- accurately enough to remove the manual-authoring cost that is the
        principal barrier to adoption --- measured both against human-judged
        developer intent and against formal discharge by an SMT-backed verifier?

  \item[RQ2 (Downstream value).] Do cheaply inferred specifications deliver
        measurable value to a demanding consumer --- specifically, supplied as
        context to a large-language-model test generator, do they improve the
        unit tests it produces relative to a signature-only baseline and a
        guidance-only control that isolates the contribution of specifications
        in particular --- so that the cost reduction is worth having?

  \item[RQ3 (Library-boundary availability).] Can inferred JML contracts be
        carried with compiled Java libraries and service-interface descriptions,
        reproduced losslessly by a consumer that has no access to the source, at
        an overhead and throughput compatible with routine workflows --- so that
        a dependency's contract is available where verification of client code
        against it must happen?

  \item[RQ4 (Compositional propagation).] Because a caller's specification is
        derived from its callees', how substantially do library specifications
        propagate into the contracts of the code that uses them? Specifically,
        does a single-pass propagation capture this cross-method dependence, or
        does a compositional analysis recover a substantial and structurally
        characterisable body of caller obligations the single pass cannot express
        --- the obligations through which a library change reaches its callers,
        and hence the basis for behavioural-compatibility analysis across
        versions?
\end{description}

\section{Contributions}

In answering these questions, the thesis makes the following contributions.

\begin{enumerate}
  \item \textbf{JML-Inferrer, a heuristic specification inference system
        organised by the WP/SP calculi.} The thesis presents the design and
        implementation of a static analyser that infers JML preconditions,
        postconditions, loop invariants, and frame conditions from Java
        source. The system's distinguishing design choice is to use the
        weakest-precondition and strongest-postcondition calculi not as
        algorithms to execute but as a \emph{discipline} that fixes which
        syntactic patterns the engine recognises and what clause each
        produces. Every emitted clause is discharged through OpenJML's
        Extended Static Checker, recovering soundness \emph{a posteriori} that
        the heuristic front end does not provide by construction.
        (Chapter~\ref{ch:inference}.)

  \item \textbf{An empirical answer, with effect size, to whether
        specifications improve LLM test generation.} A controlled experiment
        on 312 methods of Apache Commons Lang, holding the model fixed and
        varying only its input across four conditions, shows that
        specification-guided generation produces 272\% more tests and a
        mutation score 40.7 percentage points higher than signature-only
        generation. A guidance-only control rules out the alternative
        hypothesis that any prompt enrichment, not specifications in
        particular, would suffice. (Chapter~\ref{ch:testgen}.)

  \item \textbf{A toolchain-transparent mechanism for distributing inferred
        specifications.} An ASM-based embedder writes JML contracts into Java
        class-file attributes via a runtime-retained annotation, and a
        parallel OpenAPI extension family carries them across REST service
        boundaries. Evaluated on four corpora exceeding twenty thousand
        methods, the embedder achieves a 100\% lossless-roundtrip rate at a
        bytecode-size overhead below 11\%, with throughput high enough for
        per-pull-request workflows. (Chapter~\ref{ch:distribution}.)

  \item \textbf{A measurement of how strongly caller specifications depend on
        callee specifications, and its consequence for compatibility.} A
        top-down weakest-precondition refinement pass, measured across five
        open-source libraries and 21{,}052 methods, shows that a caller's
        specification depends substantially on its callees': it recovers
        42{,}370 cross-method obligations that single-pass propagation cannot
        express, dominated by the guard-conditional and dispatch-dependent
        shapes through which a library change reaches the code that calls it.
        This dependence is the mechanism that makes specifications at the
        library boundary a basis for behavioural-compatibility analysis ---
        when a library's inferred contracts change, the impact on its callers'
        contracts is computable through exactly these obligations ---
        positioning version-level compatibility checking as the capability the
        mechanism enables. (Chapter~\ref{ch:compositional}.)

  \item \textbf{An adoption argument grounded in measurement.} The thesis
        connects the four results into a single argument about why formal
        methods stay niche and what would change it: the manual cost of
        specifications is the binding constraint, inference removes it, and
        library-boundary specifications offer a capability --- verification
        across dependencies and behavioural-compatibility analysis without a
        dependency's source --- that manual specification never made routine.
        (Chapters~\ref{ch:introduction}, \ref{ch:discussion},
        and~\ref{ch:conclusion}.)

  \item \textbf{Replication artefacts.} Each study is accompanied by a
        replication package containing source code, datasets, inferred
        specifications, generated test suites, and analysis scripts,
        supporting independent verification of every empirical claim.
\end{enumerate}

\section{Research Approach and Validation Strategy}

The four research questions are answered by four empirical studies, and the
overall research approach is worth setting out, because the studies share a
methodological shape that is itself part of the thesis's argument. Each study
follows the same arc: it identifies a question that can be answered by measurement
on real code, it builds or extends the software artefact needed to take the
measurement, it runs the measurement on production libraries rather than
constructed examples, and it reports the result with its uncertainty and its
threats. The approach is empirical software engineering in the standard sense, and
its commitments --- controlled comparison, statistical treatment of stochastic
components, explicit threats to validity, replication packages --- are maintained
throughout.

The validation strategy varies with the kind of claim. For the fidelity claim
(RQ1), the validation is dual: a human oracle (two raters comparing inferred
clauses against documented intent) and a machine oracle (a verifier discharging
clauses against the implementation), reported separately because they measure
different properties. For the value claim (RQ2), the validation is a controlled
experiment with a baseline and a control condition, analysed statistically with
effect sizes and corrected significance thresholds, because the claim is causal ---
that specifications \emph{cause} better tests --- and a causal claim requires a
controlled comparison. For the reach claim (RQ3), the validation is a deterministic
measurement of losslessness, overhead, and throughput, because the claim is about
a deterministic mechanism whose properties can be measured exactly rather than
estimated statistically. For the completeness claim (RQ4), the validation is again
a deterministic measurement --- counting the clauses a second pass adds --- because
the claim is structural, supplemented by a preliminary downstream probe that
connects the structural result back to the value metric of RQ2.

The variation in validation strategy is principled: a causal claim gets a
controlled experiment, a deterministic-mechanism claim gets an exact measurement,
a structural claim gets a count. Matching the validation to the claim is what
allows each study to make a defensible assertion of the right kind --- the value
study asserts a causal effect with a confidence interval, the reach study asserts
an exact losslessness rate, and neither overclaims the other's kind of certainty.
This matching, and the consistent reporting of threats that bound each claim's
scope, is the methodological backbone the four studies share.

\section{A Note on Methodology and Stance}

Two methodological commitments run through the thesis and are worth stating
at the outset. The first is the systematic separation of \emph{heuristic
plausibility} from \emph{formal soundness}. The inference engine is
heuristic: it matches patterns, and a pattern can fire where it should not.
Rather than claim soundness the engine cannot deliver, the thesis treats
inference and verification as two independent layers. Inference proposes;
verification disposes. A clause that an OpenJML discharge rejects is flagged,
not silently retained. This stance --- infer aggressively, verify
conservatively --- recurs in every chapter and is, in a sense, the
methodological thesis underlying the technical one.

The second commitment is to measurement on real code at scale, in preference
to argument from construction. Where a claim can be made by counting --- how
many clauses discharge, how many tests a model generates, how many bytes an
embedding costs, how many clauses a second pass adds --- the thesis makes it
by counting, on production libraries (Apache Commons Lang, Commons IO,
Commons Math, jOOL, Vavr, Google Guava) rather than on toy examples. The cost
of this commitment is that the conclusions are quantitative and corpus-bound
rather than universal; the benefit is that they are defensible.

\section{The Contributions in Relation to the Literature}

The contributions can be located against the existing literature, which sharpens
what each adds. The inference engine joins a tradition that includes dynamic
detection (Daikon), verification-driven template inference (Houdini),
documentation extraction (Jdoctor), and recent language-model approaches, and it
occupies a position none of these does: purely static, so it needs no test suite;
grounded in code structure, so it needs no documentation; deterministic, so it is
consistent across runs; and validated against a formal oracle, so its heuristic
output is filtered for soundness. The engine does not dominate the alternatives on
every axis --- documentation extraction recovers exceptional behaviour the engine
misses, and dynamic detection recovers numeric invariants the engine cannot see ---
but it occupies the static-structural-verified corner that the others leave open,
and the thesis quantifies the complementarity rather than claiming dominance.

The test-generation result joins the literature on language-model test generation
and the oracle problem, where prior work has documented the problem and learned
oracles from corpora. The contribution is to show, with a controlled experiment and
an effect size, that grounded specifications supplied as context address the oracle
problem from a different direction --- not by learning oracles but by inferring
specifications --- and that the effect is large and specific to specifications
rather than to prompt enrichment in general. The guidance control is the
methodological contribution: it isolates the specification as the active
ingredient, which prior work supplying various context to models has not always
done.

The distribution mechanism addresses a gap the inference and verification
literature has largely left implicit: that a specification computed from source is
unavailable to consumers who lack the source. The contribution is a
toolchain-transparent mechanism for closing the gap, validated empirically, which
to the author's knowledge no prior work provides for JML at the bytecode level.
The compositional refinement joins the interprocedural-analysis literature, and its
contribution is a measurement --- of how substantially a caller's specification
depends on its callees' --- that the literature has not previously taken on real
code at this scale, together with the observation that this dependence is what
makes library-boundary specifications a basis for behavioural-compatibility
analysis. In each case the contribution is positioned not as a wholesale
replacement of prior work but as the filling of a specific, identifiable gap, and
the thesis's value is the demonstration that the gap-fillings together remove a
real barrier to the adoption of formal methods and open a capability ---
compatibility analysis across dependencies --- that the manual-specification regime
never made routine.

\section{Scope and Delimitations}

The thesis's claims are bounded, and stating the bounds at the outset prevents the
results from being read as more than they are. The work is confined to Java: the
inference engine, the verifier integration, and the distribution mechanism are all
Java-specific, and although the underlying ideas are not intrinsically Java-bound,
the thesis does not establish that they transfer to other languages. The evaluation
corpora are open-source library code --- utility classes, value types, collections,
numerical and functional libraries --- and the thesis does not establish its claims
on application code, framework code, concurrent code, or proprietary systems, each
of which has structural features the corpora lack. The language-model experiments
use a single model family and defend the ordering of conditions rather than the
absolute magnitudes, which the thesis does not establish are model-independent.

The thesis is also delimited in what it attempts. It does not aim to automate
verification --- it automates the specification step that verification requires as
input, leaving the verification itself to existing tools. It does not aim to produce
complete specifications --- it produces the specifications that local structural
inference can reach, and it measures the gap to documented intent rather than
closing it. It does not aim to replace the consumers it serves --- it supplies
specifications to test generators and verifiers rather than competing with them.
And it does not aim to settle the design question of whether a second compositional
pass should always run --- it measures the gap the single pass leaves and provides
the information a deployer needs to decide, leaving the decision to the deployment
context. These delimitations are not weaknesses to be apologised for but the
boundaries that make the claims precise: the thesis establishes specific results on
specific corpora with specific tools, and the value of the results is in their
defensibility within those bounds, not in an overreach beyond them.

\section{Thesis Structure}

The remainder of this thesis is organised as follows.

\textbf{Chapter~\ref{ch:background}} surveys the background and related work
spanning the four research questions: the program-reasoning calculi (Hoare
logic, weakest preconditions, strongest postconditions) and the Java Modeling
Language that operationalises them; the landscape of specification inference,
from dynamic detection (Daikon) through documentation-based extraction
(Jdoctor) to recent language-model approaches; the literature on
language-model test generation and the oracle problem; mutation testing as an
evaluation methodology; the verification tools that serve as the formal
oracle; and the channels by which specifications might be distributed.

\textbf{Chapter~\ref{ch:inference}} presents JML-Inferrer: its architecture,
its WP/SP-organised inference of preconditions, postconditions, loop
invariants, and interprocedural summaries, its OpenJML-backed validation
layer, and its method categorisation scheme.

\textbf{Chapter~\ref{ch:testgen}} reports the controlled experiment on
specification-guided LLM test generation, including the manual and formal
validation of the inferred specifications that the experiment depends on.

\textbf{Chapter~\ref{ch:distribution}} develops the specification-distribution
mechanism --- the class-file embedding and the OpenAPI extension --- and
reports its empirical validation across four corpora.

\textbf{Chapter~\ref{ch:compositional}} measures how substantially a caller's
specification depends on its callees', via a top-down compositional analysis
evaluated on five libraries, and develops the consequence of that dependence ---
that a library change propagates into its callers' contracts --- as the basis for
behavioural-compatibility analysis at the library boundary. It also reports a
preliminary downstream signal from re-running the test-generation experiment with
the more thoroughly propagated specifications.

\textbf{Chapter~\ref{ch:discussion}} synthesises the four studies around the
adoption question, drawing out the cross-cutting themes --- the cost barrier and
its removal, the library-boundary capabilities inferred specifications enable, the
recurring heuristic-versus-formal tension, and the combined threats to validity
--- and \textbf{Chapter~\ref{ch:conclusion}} concludes and sets out directions for
future work, the foremost of which is the full version-to-version compatibility
study the compositional mechanism makes possible.

The chapters are meant to be read in sequence, because the argument is cumulative:
inference produces the specifications (Chapter~\ref{ch:inference}), the
test-generation study shows they are worth having (Chapter~\ref{ch:testgen}),
distribution carries them to the library boundary (Chapter~\ref{ch:distribution}),
and the compositional study shows how they propagate across that boundary and what
that propagation enables (Chapter~\ref{ch:compositional}), before
Chapter~\ref{ch:discussion} draws the threads into the adoption argument. Each
chapter builds on the background and methods established once in
Chapters~\ref{ch:background} and~\ref{ch:inference} rather than restating them, and
each hands off explicitly to the next. What unifies them is the single question
the thesis pursues throughout: why formal methods stay niche, and whether removing
the specification-authoring cost and making specifications a property of the
library boundary can change that.
